% generated by extract_options.py from src/search/sat/bdd_encoding.cc -- do not edit
\begin{longtable}{@{}p{0.27\linewidth}p{0.13\linewidth}p{0.54\linewidth}@{}}
\toprule
option & default & meaning \\
\midrule
\endhead
\texttt{one\_step\_only} & \texttt{true} & only allow one real transition per factor and time step (plus self loops before and after it) instead of arbitrary label sequences \\
\texttt{combinebdds} & \texttt{false} & use one BDD per factor over (state, next state, labels) instead of one BDD per (factor, source, target) triple conditioned on the state variables \\
\texttt{bdd\_size\_limit} & \texttt{-1} & per-factor budget of BDD nodes. Oversized (source,target) BDDs fall back to the one-step BDD. -1 means no limit. Requires one\_step\_only=false \\
\texttt{impltseitsin} & \texttt{true} & use implicational (one-directional) Tseitin encoding instead of a bi-implicational one \\
\texttt{path\_encoding} & \texttt{false} & translate each BDD with a path-based encoding (NNFPath, Abio et al. CPAIOR 2016) instead of Tseitin: variables for the nodes and edges on the root-to-1 path, plus clauses that make unit propagation domain consistent on the BDD's variables. About three variables per node and no sharing between BDDs; best with combinebdds \\
\texttt{path\_sharing} & \texttt{false} & with path\_encoding and split BDDs: one path encoding per factor over all its (source, target) BDDs, which share nodes; the BDD of the active pair is selected by an auxiliary variable for previous\_state(s) \& next\_state(s'). Without it every pair's BDD gets its own path encoding \\
\texttt{label\_implications\_bicliques} & \texttt{false} & encode the binary label\_implications clauses by a greedy biclique cover (Rintanen, ECAI 2006): a biclique A x B of conflicting literals becomes a -\textgreater{} z, b -\textgreater{} -z with one new variable per step instead of |A|*|B| binary clauses \\
\texttt{combine\_ite} & \texttt{false} & with combinebdds: build the factor BDD as if-then-else chains over the state variables instead of conjoining full one-hot minterms. Equal on one-hot assignments, which the exactly-one constraints guarantee; |S|\^{}2 instead of |S|\^{}3 BDD operations \\
\texttt{omitforcedvariables} & \texttt{true} & do not create Tseitin variables for BDD nodes whose in-degree is at most forcedvariablesthreshold \\
\texttt{forcedvariablesthreshold} & \texttt{100} & in-degree up to which omitforcedvariables applies \\
\texttt{cutbdds} & \texttt{false} & fixpoint 'cutting' of the BDDs across factors for stronger constraints \\
\texttt{cut\_mode} & \texttt{AND} & with cutbdds: how a constraint is applied \emph{(AND, AND\_LIMITED, RESTRICT, GLOBAL\_AND, GLOBAL\_AND\_LIMITED, GLOBAL\_RESTRICT, AND\_CAPPED)} \\
\texttt{cut\_limit\_factor} & \texttt{1} & with cut\_mode=AND\_LIMITED \\
\texttt{report\_label\_conjunction} & \texttt{false} & diagnostics: conjoin the any-transition label BDDs of all factors (every label set permitted by every factor in one step) under label\_conjunction\_limit new nodes per step, and report size, time and how many factors fit \\
\texttt{label\_conjunction\_limit} & \texttt{10000000} & node limit for report\_label\_conjunction \\
\texttt{label\_implications} & \texttt{NONE} & add clauses implied by the factor BDDs, so unit propagation sees dependencies the Tseitin translation does not propagate \emph{(NONE, UNCONDITIONAL, PER\_SOURCE, BOTH)} \\
\texttt{label\_implications\_max} & \texttt{-1} & at most this many implied clauses of each kind (unconditional, per source); -1 = all. Extraction stops as soon as the cap is reached; the clauses found so far are kept \\
\texttt{label\_implications\_time\_limit} & \texttt{-1} & stop extracting implied clauses after this many seconds and keep the ones found so far (all of them are sound); -1 = no limit \\
\texttt{source\_units} & \texttt{false} & split BDDs: prev\_s -\textgreater{} +-l for the labels forced in every label set the factor permits from s \\
\texttt{target\_units} & \texttt{false} & split BDDs: next\_s' -\textgreater{} +-l for the labels forced in every label set leading into s' \\
\texttt{pair\_units} & \texttt{false} & split BDDs: prev\_s \& next\_s' -\textgreater{} +-l for the labels the (s, s') BDD forces (not already source/target units) \\
\texttt{cross\_units} & \texttt{false} & split BDDs: for two factors sharing labels, prev\_s1 \& prev\_s2 -\textgreater{} +-l (and the same for next states) for labels forced only by both together; factor pairs with at most cross\_pair\_limit state pairs, most shared labels first \\
\texttt{cross\_pair\_limit} & \texttt{10000} & cross\_units: state pairs per factor pair \\
\texttt{leave\_clauses} & \texttt{false} & prev\_s \& -next\_s -\textgreater{} some label leaving s fires; prev\_s \& l -\textgreater{} a label leaving s fires before l (in the label order) for every label l without a transition from s \\
\texttt{enter\_clauses} & \texttt{false} & next\_s' \& -prev\_s' -\textgreater{} some label entering s' fires; next\_s' \& l -\textgreater{} a label entering s' fires after l for every label l without a transition into s' \\
\texttt{state\_clauses\_time\_limit} & \texttt{30} & overall budget for the state clauses, computed in phases leave/enter, source/target units, pair units, cross units; what a phase found before it stopped is kept; -1 = no limit \\
\texttt{state\_clauses\_min\_phase} & \texttt{2} & seconds of the budget reserved for each enabled phase still to come: phase k stops at budget - min\_phase * (enabled phases after k) \\
\texttt{move\_clause\_parts} & \texttt{BOTH} & leave/enter clauses: which of the two parts to emit \emph{(BOTH, FRAME, ORDER)} \\
\texttt{move\_clauses\_max} & \texttt{1000000} & leave/enter clauses: at most this many (state, blocked label) entries; -1 = all \\
\texttt{goal\_last\_label} & \texttt{false} & at the last step, a label that cannot move any factor into its goal states must be followed by a later label (a plan whose last action achieves no goal fact can be shortened) \\
\texttt{cut\_min\_overlap} & \texttt{0} & cutbdds (pairwise modes): only propagate from S to T if the labels they share make up at least this fraction of T's labels \\
\texttt{cube\_encoding\_max} & \texttt{-1} & encode a BDD with at most this many paths to 0 directly, one clause per path (cube of its negation) together with the conditions, without Tseitin variables; larger BDDs keep the Tseitin encoding. -1 = off (Tseitin only) \\
\texttt{coverbdds} & \texttt{false} & search for always-true/false state -\textgreater{} +-label implications and report them. Diagnostic only: this does not change the generated formula \\
\texttt{force\_at\_least\_one\_action} & \texttt{false} & force that every time step contains at least one action \\
\texttt{alternate\_label\_orders} & \texttt{false} & use two label orderings instead of one: the order from label\_order for the odd-numbered time steps and its reverse for the even ones. The label order is the BDD variable order, so this means building and encoding two sets of BDDs, and a step's label variables are then constrained in the order belonging to that step \\
\texttt{report\_label\_implications} & \texttt{false} & mine each factor's label BDD for dependencies that hold in every legal label set (l, -l, l -\textgreater{} l', l -\textgreater{} -l') and report how many there are. Diagnostic only: it does not change the formula \\
\texttt{report\_factor\_statistics} & \texttt{false} & emit one BDDSTAT line per factor with its sizes and construction time. Off by default: it is one line per factor and there can be hundreds. The single construction\_ok/construction\_failed summary is always emitted \\
\texttt{cudd\_init\_nodes} & \texttt{16000000} & initial size of the Cudd unique table, divided by the number of labels \\
\texttt{cudd\_cache\_size} & \texttt{16000000} & initial number of Cudd cache slots. Allocated up front and independent of task size, so it dominates memory on small tasks \\
\texttt{cudd\_max\_memory\_mb} & \texttt{0} & hard cap on Cudd's memory in MB. 0 lets Cudd decide from the machine's RAM, which ignores any per-run budget \\
\texttt{bdd\_init\_time\_limit} & \texttt{-1} & seconds allowed for building the BDDs. When it is used up the run stops with 'BDDSTAT construction\_failed ... reason time\_limit' instead of grinding on until the driver kills it. -1 means no limit \\
\texttt{bdd\_node\_limit} & \texttt{-1} & how many live Cudd nodes the construction may use, reported the same way. -1 means no limit \\
\texttt{label\_order} & \texttt{label\_order\_linear()} & order in which labels may be applied within one time step. Unless bdd\_variable\_order is given, this is also the BDD variable order \\
\texttt{bdd\_variable\_order} & \texttt{} & order of the label variables in the BDDs, independently of label\_order. It changes only the size of the BDDs, not which label sets they allow. Omitted: the label order \\
\texttt{bdd\_reorder} & \texttt{NONE} & Cudd variable reordering, applied once after the BDDs are built (and after cutbdds), before they are Tseitin-translated. The labels keep their meaning; only the variable levels change \emph{(NONE, SIFT, SIFT\_CONVERGE, SYMM\_SIFT, GROUP\_SIFT)} \\
\texttt{bdd\_reorder\_dynamic} & \texttt{false} & also let Cudd reorder automatically during construction, with the bdd\_reorder method \\
\texttt{dump\_bdd\_variable\_order} & \texttt{none} & write the final BDD variable order of the labels (after reordering, top level first) to this file, one label id per line, in the format label\_order\_file reads; 'none' = do not. The option parser lower-cases the path. \\
\texttt{dump\_label\_order} & \texttt{none} & write the (primary) label order to this file, one label id per line; 'none' = do not. The option parser lower-cases the path. \\
\bottomrule
\end{longtable}
